% =========================================================================
% Préambule du mémoire
% =========================================================================

% --- Encodage et langue ---------------------------------------------------
\usepackage[utf8]{inputenc}
\usepackage[T1]{fontenc}
\usepackage[french]{babel}

% --- Mise en page ---------------------------------------------------------
\usepackage[a4paper, margin=2.5cm]{geometry}
\usepackage{microtype}
\usepackage{setspace}
\onehalfspacing

% --- Couleurs et liens ----------------------------------------------------
\usepackage{xcolor}
\definecolor{linkcolor}{RGB}{20, 60, 120}
\definecolor{codebg}{RGB}{248, 248, 248}
\definecolor{codeframe}{RGB}{200, 200, 200}

\usepackage[
  colorlinks=true,
  linkcolor=linkcolor,
  citecolor=linkcolor,
  urlcolor=linkcolor,
  pdftitle={Simulation comportementale d'un aquarium en C/GTK4 avec configuration externe},
  pdfauthor={Yahya},
  pdfsubject={Projet PFA},
  bookmarksopen=true,
  bookmarksnumbered=true
]{hyperref}

% --- Graphiques et tableaux -----------------------------------------------
\usepackage{graphicx}
\usepackage{caption}
\usepackage{subcaption}
\usepackage{booktabs}
\usepackage{array}
\usepackage{longtable}
\usepackage{enumitem}
\usepackage{float}

% --- Mathématiques --------------------------------------------------------
\usepackage{amsmath}
\usepackage{amssymb}

% --- TikZ -----------------------------------------------------------------
\usepackage{tikz}
\usetikzlibrary{
  arrows.meta,
  positioning,
  shapes,
  shapes.geometric,
  automata,
  calc,
  fit,
  decorations.pathreplacing,
  backgrounds
}

% --- Code source via minted ------------------------------------------------
\usepackage[newfloat]{minted}
\setminted{
  fontsize=\footnotesize,
  linenos,
  numbersep=8pt,
  frame=lines,
  framesep=2mm,
  breaklines,
  breakanywhere=false,
  bgcolor=codebg,
  tabsize=4
}

% Style par défaut pour le C
\setminted[c]{style=friendly}
\setminted[toml]{style=friendly}
\setminted[bash]{style=friendly}
\setminted[make]{style=friendly}

% Liste flottante pour les codes (« Liste des codes »).
% Avec l'option [newfloat], minted utilise le package newfloat qui
% n'expose pas \listingscaption ni \listoflistingscaption : on les
% renomme via caption et on redéfinit la macro pour la liste.
\captionsetup[listing]{position=above, name=Listing}
\providecommand{\listoflistingscaption}{Liste des codes}
\renewcommand{\listoflistingscaption}{Liste des codes}
\providecommand{\listoflistings}{\listof{listing}{\listoflistingscaption}}

% csquotes doit être chargé après fvextra (chargé par minted) pour éviter
% un warning de lineno.
\usepackage{csquotes}

% --- Bibliographie (BibTeX classique, style plain) -----------------------
% biblatex n'est pas disponible sur cette installation : on utilise
% le système traditionnel \bibliography{...} + \bibliographystyle{...}.

% --- Numérotation des pages -----------------------------------------------
% Numérotation arabe (1, 2, 3, ...) sur tout le document.
\pagenumbering{arabic}

% --- Profondeur de la table des matières ----------------------------------
\setcounter{tocdepth}{2}
\setcounter{secnumdepth}{3}

% --- Commandes personnalisées ---------------------------------------------
\newcommand{\code}[1]{\texttt{#1}}
\newcommand{\file}[1]{\texttt{#1}}
\newcommand{\struct}[1]{\texttt{#1}}
\newcommand{\func}[1]{\texttt{#1()}}
\newcommand{\kw}[1]{\textbf{\texttt{#1}}}

% Mise en avant d'un terme défini
\newcommand{\term}[1]{\emph{#1}}

% --- En-têtes et pieds de page --------------------------------------------
\usepackage{fancyhdr}
\setlength{\headheight}{14pt}
\addtolength{\topmargin}{-2pt}
\pagestyle{fancy}
\fancyhf{}
\fancyhead[L]{\leftmark}
\fancyhead[R]{\thepage}
\fancyfoot[C]{}
\renewcommand{\headrulewidth}{0.4pt}
\renewcommand{\footrulewidth}{0pt}

% Empêche le mode "lettre majuscule partout" dans \leftmark
\renewcommand{\chaptermark}[1]{\markboth{\chaptername\ \thechapter\ -- #1}{}}

% --- Espacement avant les chapitres ---------------------------------------
\usepackage{titlesec}
\titleformat{\chapter}[display]
  {\normalfont\huge\bfseries}
  {\chaptertitlename\ \thechapter}
  {20pt}
  {\Huge}
\titlespacing*{\chapter}{0pt}{-30pt}{20pt}
